\chapter{Language essentials and HTTP boundaries}

The A-to-Z example already showed the main building blocks in action. This chapter turns those pieces into a day-to-day development model for web work: organize source explicitly, terminate statements unambiguously, compute with typed values, normalize text at boundaries, keep domain logic in handlers, and cross into HTTP/HTML only through typed routes and safe rendering.

A useful working rule is: \emph{parse and normalize at the boundary, compute in typed values, persist through explicit queries, and render as late as possible}. The sections below follow that order.

\section{Source structure and statement boundaries}

\subsection*{Modules}

\code{main.rw} is the entry point of the module graph:

\begin{lstlisting}[language=RWLang]
mod models;
mod queries;
mod components;
mod pages;
mod actions;
\end{lstlisting}

In V1, \code{mod} loads an application-root-relative module without injecting its declarations into the global namespace. For example, declarations from \code{models.rw} are addressed as \code{models::Article} across module boundaries. Local declarations may use their short name inside their own module. There is no general Rust-like \code{pub} export system in V1.

\subsection*{Module namespaces and cross-module calls}

A module declaration adds one application-root-relative source module to the explicit source graph; it does not copy that module's declarations into the current scope. The canonical mapping is direct:

\begin{lstlisting}
models.rw              -> models
catalog/queries.rw     -> catalog::queries
catalog/pages.rw       -> catalog::pages
\end{lstlisting}

For example, an entry point may load several modules:

\begin{lstlisting}[language=RWLang]
mod models;
mod catalog::queries;
mod catalog::pages;

route catalogIndex GET "/catalog"
    => catalog::pages::index;
\end{lstlisting}

Across a module boundary, use the qualified name. A declaration in \code{catalog/queries.rw} is referenced as \code{catalog::queries::recent(...)}, and a type in \code{models.rw} as \code{models::Article}. Inside the declaration's own module, the short local name remains valid. This keeps local code concise while making dependencies between source modules visible.

Module paths are always relative to the application root. There is no \code{mod.rw} fallback, automatic directory discovery, \code{../}, \code{./}, \code{self::}, \code{super::}, \code{crate::}, or wildcard import. Declare every source module explicitly, before ordinary top-level declarations.

Language builtins such as \code{trim(...)}, \code{sin(...)}, and \code{len(...)} are not application-module declarations and therefore remain unqualified.

\subsection*{Statement terminators}

RWLang uses an explicit semicolon as the terminator for simple statements. Newlines are whitespace, not statement boundaries; there is no automatic semicolon insertion. This keeps multiline expressions deterministic and makes generated/audited code unambiguous.

\begin{lstlisting}[language=RWLang]
let subtotal = price
    * quantity
    + shipping;
set retries = retries + 1;
authorize article owner authorUsername or role Publisher;
return Ok(json(subtotal));
\end{lstlisting}

Block statements and block declarations are closed by their brace and do not take a trailing semicolon:

\begin{lstlisting}[language=RWLang]
if ready {
    set state = 1;
}
while state < 3 {
    set state = state + 1;
}
\end{lstlisting}

The same rule applies inside a transaction: a standalone mutating query call and a business-audit statement are simple statements and therefore end in \code{;}. Top-level block declarations such as \code{model}, \code{page fn}, \code{action fn}, and control-flow blocks do not use \code{;}. Non-block top-level declarations do: both \code{mod path;} and \code{route ... => handler;} are terminated explicitly with \code{;}. A multiline route is complete only at that final terminator.

\section{Typed computation for web application code}

\subsection*{Expressions and arithmetic}

RWLang expressions use explicit precedence. From tighter to looser: unary \code{!}; multiplication, division and remainder; addition and subtraction; shifts; comparisons; bitwise AND/XOR/OR; logical AND; logical OR. Parentheses should be used whenever they make intent clearer.

\begin{lstlisting}[language=RWLang]
let total = 12 + 3 * 4;
let remainder = 17 % 5;
let shifted = 3 << 2;
let masked = 12 & 10;
let allowed = ready && authorized;
let fallback = cached || fresh;
let denied = !allowed;
\end{lstlisting}

The operator contract is deliberately typed:

\begin{center}
\begin{tabularx}{\textwidth}{@{}l l Y@{}}
\toprule
\textbf{Operands} & \textbf{Operators} & \textbf{Result} \\
\midrule
\code{Int}, same type & \code{+ - * / \%} & checked \code{Int}. \\
\code{F32}, same type & \code{+ - * / \%} & \code{F32}; result must remain finite. \\
\code{Decimal}, same type & \code{+ - * / \%} & checked \code{Decimal}. \\
\code{Int}, same type & \code{<< >> \& \^{} |} & shift/bitwise \code{Int}. \\
\code{Bool}, same type & \code{\&\& ||} & \code{Bool}; short-circuit evaluation. \\
\code{Bool} & \code{!} & logical negation. \\
\code{String}, same type & \code{+} & \code{String} concatenation, allocation-budgeted. \\
Numeric, same type & \code{< <= > >=} & \code{Bool}. \\
Same supported scalar type & \code{== !=} & \code{Bool}. \\
\bottomrule
\end{tabularx}
\end{center}

There is no implicit mixing of \code{Int}, \code{F32}, and \code{Decimal}. Convert an integer explicitly with \code{toF32(value)} when floating-point computation is required. Numeric overflow, division/remainder by zero, invalid shifts, or a non-finite \code{F32} result fail closed rather than silently wrapping or propagating NaN/Infinity.

The numeric builtins include:

\begin{center}
\begin{tabularx}{\textwidth}{@{}l l Y@{}}
\toprule
\textbf{Function} & \textbf{Type} & \textbf{Purpose} \\
\midrule
\code{sin(x)}, \code{cos(x)} & \code{F32 -> F32} & Trigonometric functions. \\
\code{sqrt(x)} & \code{F32 -> F32} & Square root; invalid domain fails. \\
\code{abs(x)} & \code{Int -> Int} or \code{F32 -> F32} & Absolute value. \\
\code{ln(x)}, \code{log10(x)} & \code{F32 -> F32} & Natural/base-10 logarithm. \\
\code{log(x, base)} & \code{F32, F32 -> F32} & Logarithm with explicit base. \\
\code{exp(x)} & \code{F32 -> F32} & Exponential function. \\
\code{pow(x, y)} & \code{F32, F32 -> F32} & Floating-point power. \\
\code{round/floor/ceil(x)} & \code{F32 -> F32} & Rounding operations. \\
\code{toF32(x)} & \code{Int -> F32} & Explicit integer-to-\code{F32} conversion. \\
\code{monotonicNanos()} & \code{() -> Int} & Process-local monotonic nanosecond counter. \\
\bottomrule
\end{tabularx}
\end{center}

All \code{F32} builtin results must remain finite. Invalid-domain operations fail instead of creating NaN or infinity.\par
\code{monotonicNanos()} is for elapsed-time measurement, not wall-clock timestamps; public-cache output that depends on it is rejected.

\subsection*{Practical arithmetic patterns}

Web code rarely needs advanced mathematics everywhere, but a few numeric patterns recur constantly. Keep them typed and visible rather than hiding them in string conversions or SQL fragments.

For pagination, compute an offset from validated integers:

\begin{lstlisting}[language=RWLang]
let pageSize = 25;
let offset = (page - 1) * pageSize;
\end{lstlisting}

For flags and compact integer masks, bitwise operations are available, but they should represent a real integer-domain contract rather than replace readable booleans. Use \code{&&} and \code{||} for business conditions because they short-circuit:

\begin{lstlisting}[language=RWLang]
let mayPublish = authenticated
    && ownsArticle
    && !locked;
\end{lstlisting}

Use \code{Decimal} for money. Use \code{F32} only when floating-point behavior is part of the problem, such as geometry, signal processing, statistics, or explicit math functions. Never convert monetary values to \code{F32} merely to reuse a math builtin.

\subsection*{String and text functions}

String operations are typed, Unicode-aware builtins executed inside the bounded runtime:

\begin{lstlisting}[language=RWLang]
let cleaned = trim(text);
let left = trimStart(text);
let right = trimEnd(text);
let normalized = lower(cleaned);
let chars = stringLen(cleaned);
let rewritten = replace(normalized, " ", "-");
let part = substring(cleaned, 1, 3);
let first = indexOf(cleaned, "rw");
let last = lastIndexOf(cleaned, "a");
let ch = charAt("RWLang", 1);
let repeated = repeat("rw", 3);
let parts = split(rewritten, "-");
\end{lstlisting}

\begin{center}
\begin{tabularx}{\textwidth}{@{}Y Y@{}}
\toprule
\textbf{Function} & \textbf{Signature} \\
\midrule
\code{stringLen(text)} & \code{String -> Int} \\
\code{trim/trimStart/trimEnd(text)} & \code{String -> String} \\
\code{lower/upper(text)} & \code{String -> String} \\
\code{contains(text, needle)} & \code{String, String -> Bool} \\
\code{startsWith/endsWith(text, part)} & \code{String, String -> Bool} \\
\code{replace(text, search, replacement)} & \code{String, String, String -> String} \\
\code{split(text, delimiter)} & \code{String, String -> List<String>} \\
\code{substring(text, start[, length])} & \code{String, Int[, Int] -> String} \\
\code{indexOf/lastIndexOf(text, needle)} & \code{String, String -> Int} \\
\code{charAt(text, index)} & \code{String, Int -> String} \\
\code{repeat(text, count)} & \code{String, Int -> String} \\
\bottomrule
\end{tabularx}
\end{center}

\code{stringLen}, \code{substring}, \code{indexOf}, \code{lastIndexOf}, and \code{charAt} use Unicode scalar-value positions rather than UTF-8 byte offsets. Missing searches return \code{-1}; invalid negative/out-of-range indices fail closed. \code{replace} rejects an empty search string, \code{split} rejects an empty delimiter and is capped at 4096 items, while \code{repeat} is allocation-budgeted before materializing its result. \code{List<String>} remains a typed compute collection with \code{len(parts)} and checked \code{parts[index]} access.

For bounded regular-expression work, RWLang also provides \code{regexMatch}, \code{regexReplace}, and \code{regexCaptures}. Regex patterns and runtime work remain subject to validation and resource limits.

\subsection*{Practical text-handling patterns}

Treat text manipulation as boundary work. Normalize user-facing search/filter input before comparing it, but keep the original value when the exact spelling matters to the domain.

\begin{lstlisting}[language=RWLang]
let queryText = trim(q);
let normalizedQuery = lower(queryText);
let hasRwlang = contains(normalizedQuery, "rwlang");
\end{lstlisting}

Use \code{replace} for deterministic substitutions, not for general parsing:

\begin{lstlisting}[language=RWLang]
let displayTag = replace(trim(tag), "_", " ");
\end{lstlisting}

When slicing or locating text, remember that RWLang uses Unicode scalar-value positions rather than UTF-8 byte offsets. This makes \code{substring}, \code{charAt}, and the index functions suitable for application text without exposing byte-level encoding details. For structured identifiers such as URLs, slugs, emails, dates, or UUIDs, prefer the dedicated RWLang types and validators instead of rebuilding their rules with string functions.

\subsection*{Choosing the right tool in application code}

\begin{center}
\begin{tabularx}{\textwidth}{@{}l Y@{}}
\toprule
\textbf{Web-development task} & \textbf{Preferred RWLang approach} \\
\midrule
Normalize search/filter text & \code{trim}, then \code{lower}/\code{upper} only when comparison should ignore case. \\
Replace a known token & \code{replace}; do not use regex unless pattern matching is actually required. \\
Parse email, URL, slug, date, UUID & Use the dedicated typed value and validation path, not ad-hoc string parsing. \\
Money/totals & \code{Decimal}. \\
Pagination/counters & \code{Int} with checked arithmetic. \\
Business conditions & \code{Bool} with short-circuit \code{&&}, \code{||}, and \code{!}. \\
Scientific/statistical math & \code{F32} and explicit math builtins; check that floating point is appropriate for the domain. \\
Elapsed-time measurement & \code{monotonicNanos()}, never as a wall-clock timestamp. \\
\bottomrule
\end{tabularx}
\end{center}

\section{Domain values and handler code}

\subsection*{Models and business types}

\begin{lstlisting}[language=RWLang]
model Article {
    id: Int
    slug: Slug
    title: String
    body: String
}
\end{lstlisting}

Common types include \code{String}, \code{Int}, and \code{Bool}, as well as \code{Date}, \code{DateTime}, \code{Uuid}, \code{Decimal}, and \code{Slug}. Use \code{Decimal} for money rather than floating-point values.

\subsection*{Page and action}

A page contains GET-like presentation logic:

\begin{lstlisting}[language=RWLang]
page fn show(ctx: PageContext, db: Db, id: Int)
    -> Result<Html, PageError> {
    let article = articleById(db, id)?;
    return Ok(html { <h1>{{ article.title }}</h1> });
}
\end{lstlisting}

An action handles state-changing requests:

\begin{lstlisting}[language=RWLang]
action fn create(ctx: ActionContext, db: Db, title: String)
    -> Result<Redirect, PageError> {
    transaction db {
        insertArticle(tx, title)?;
    }
    return Ok(redirect("/admin/articles"));
}
\end{lstlisting}

\subsection*{Variables and error propagation}

\code{let} introduces a typed local value. V1 has no general-purpose \code{mut}; compute code uses explicit \code{set name = expr} reassignment, with the original static type preserved and execution charged to the instruction budget. The \code{?} operator propagates an error from a typed \code{Result}.

\subsection*{Enum}

Use an enum for a closed set of choices rather than free-form text. This gives a stronger contract in forms, routes, JSON, and domain logic.

\subsection*{Object}

Related model/query/page operations can be grouped under a domain object. This is particularly useful in larger applications because it defines a clear domain API boundary.

\section{Routing, typed input, and HTML}

\subsection*{Route}

\begin{lstlisting}[language=RWLang]
route article GET "/article/:slug<Slug>" => article;
\end{lstlisting}

The path parameter type is declared in the route. Invalid input never reaches the handler as a typed value.

\subsection*{HTML escaping}

\begin{lstlisting}[language=RWLang]
return Ok(html {
    <h1>{{ article.title }}</h1>
});
\end{lstlisting}

\code{{ ... }} is HTML-escaped. A string read from the database is no exception. There is no raw-HTML API.

\subsection*{Lists and optional values}

A list must be iterated, and an optional model may be dereferenced only after an explicit presence check. The compiler does not allow a \code{List<Model>} or an unchecked \code{Model?} to behave like a single record.

\begin{lstlisting}[language=RWLang]
@for article in articles {
    <li>{{ article.title }}</li>
}

@if article {
    <h1>{{ article.title }}</h1>
}
\end{lstlisting}

The database chapter explains the exact runtime query-cardinality contract, including the fact that a \code{Model?} result is an error if the query returns more than one row.

\subsection*{Typed URLs}

\begin{lstlisting}[language=RWLang]
<a @href(article, article.slug)>Open</a>
<form method="post" @action(deleteArticle, article.id)>
\end{lstlisting}

\subsection*{Prefer route helpers}

Do not duplicate your own application-route URLs as hand-written strings. In HTML, \code{@action(...)} is the preferred form for POST targets and \code{@href(...)} for links, because the route name and typed parameters remain the source of truth. Built-in runtime endpoints -- such as \code{/__rw/auth/logout} -- are naturally not application route helpers.

Do not build a dynamic \code{href="{{ value }}"} link. A route helper knows both the destination and the parameter type.

\subsection*{Components and layouts}

\begin{lstlisting}[language=RWLang]
layout fn Main(title: String) -> Html {
    html {
        <html>
            <head><title>{{ title }}</title></head>
            <body><main>@content</main></body>
        </html>
    }
}
\end{lstlisting}

Inside a page:

\begin{lstlisting}[language=RWLang]
return Ok(html {
    @layout(Main, article.title) {
        <article><h1>{{ article.title }}</h1></article>
    }
});
\end{lstlisting}

A layout may contain exactly one \code{@content} slot; components and layouts can only access their explicit typed parameters.
